\documentclass[conference]{IEEEtran}

% -- Packages --------------------------------------------------------------
\usepackage{amsmath,amssymb,bm}
\usepackage{graphicx}
\usepackage[caption=false]{subfig}
\usepackage{booktabs}
\usepackage{multirow}
\usepackage{array}
\usepackage{xcolor}
\usepackage{hyperref}
\usepackage{cleveref}

% -- Macros ----------------------------------------------------------------
\newcommand{\States}{\mathcal{S}}
\newcommand{\Actions}{\mathcal{A}}
\newcommand{\Obs}{\mathcal{O}}
\newcommand{\Trans}{T}
\newcommand{\Rew}{R}
\newcommand{\ObsFn}{Z}
\newcommand{\Disc}{\gamma}
\newcommand{\belief}{b}
\newcommand{\pf}{\hat{b}}
\newcommand{\ee}{\mathbf{p}_{\text{ee}}}
\newcommand{\block}{\mathbf{p}_{\text{blk}}}
\newcommand{\goal}{\mathbf{p}_{\text{goal}}}
\newcommand{\sigmaobs}{\sigma_{\text{obs}}}
\newcommand{\sigmaep}{\sigma_{\text{ep}}}

% -- Title -----------------------------------------------------------------
\title{Belief-Augmented Reinforcement Learning for Pick-and-Place
       Under Observation Noise and Occlusion}

\author{
  \IEEEauthorblockN{Allen Devaraj Augustin Ponraj}
  \IEEEauthorblockA{
    Department of Aerospace Engineering Sciences\\
    University of Colorado Boulder\\
    allendevaraj.augustinponraj@colorado.edu
  }
}

\begin{document}
\maketitle

% -- Abstract --------------------------------------------------------------
\begin{abstract}
Robotic pick-and-place under partial observability is challenging because
sensor noise and occlusion corrupt or eliminate target state information,
turning a standard MDP into a POMDP.
Three solution strategies are investigated on the SO-ARM101 6-DOF arm in
a MuJoCo environment featuring distance-dependent wrist camera noise and
two occluding distractor blocks:
(i) plain PPO with a stale-sensor observation model,
(ii) belief-augmented PPO that feeds a bootstrap particle filter's mean
and uncertainty estimate directly into the policy observation,
and (iii) POMCP-MCTS online planning as a stretch goal.
The central finding is that particle filter belief augmentation provides
a statistically meaningful advantage in the moderate-noise regime,
peaking at $+24\%$ task success at $\sigmaobs = 10$~mm,
but is outperformed by plain PPO at very low noise ($\sigmaobs = 3$~mm, $-30\%$)
where raw observations are already accurate and filter initialization overhead
dominates.
Both methods degrade at extreme noise ($30$~mm) due to filter impoverishment
and stale-sensor failure respectively.
The crossover near $\sigmaobs \approx 8$~mm is theoretically coherent:
it marks the point at which filtering benefit first exceeds convergence cost,
and is consistent with the signal-to-noise characteristics of a bootstrap
particle filter with $N\!=\!300$ particles.
\end{abstract}

% -- I. Introduction -------------------------------------------------------
\section{Introduction}
\label{sec:intro}

Robotic manipulation in unstructured environments requires acting reliably
despite imperfect sensing.
Real cameras return noisy detections, distractors occlude the target, and
physical contact changes the scene in ways the robot cannot directly observe.
Treating such settings as fully observable MDPs (as most deep RL work does)
ignores the fundamental epistemic uncertainty in the system and can lead to
brittle policies that fail when sensor quality degrades or objects move
out of view.

This work examines a concrete instance of the problem: pick-and-place of a
target LEGO block in the presence of two occluding distractor blocks, using
only a wrist-mounted camera whose noise scales with distance to the target.
Three distinct uncertainty sources are present simultaneously.
\emph{Sensor noise:} the wrist camera returns a Gaussian-corrupted position
estimate whose standard deviation grows linearly with the end-effector-to-target
distance, so the agent must approach the target to improve its own observations.
\emph{Occlusion:} when a distractor block passes between the camera and the
target, the observation is entirely absent; the agent must remember and reason
about the last known position.
\emph{Distractor identity:} because the auto-grasp threshold ($10$~mm) is
comparable to the inter-block spacing ($\sim\!17$~mm mean), the agent risks
grasping the wrong block if it navigates to a stale or incorrect position estimate.
Together these make the task a POMDP rather than an MDP: the true target
pose is hidden, and observations are noisy, intermittent, and corrupted by
irrelevant objects.

Two learned policy approaches are compared.
\emph{Plain PPO} treats the problem as memoryless: it receives the last visible
wrist observation (replayed during occlusion) and learns to act on this noisy
signal directly.
\emph{Belief-augmented PPO} maintains a bootstrap particle filter over the
target's $(x, y, \theta)$ pose and augments the policy observation with the
filter's mean $\bm{\mu}$ and standard deviation $\bm{\sigma}$, giving the
policy an explicit uncertainty signal it can condition its behavior on.
Both policies are trained with proximal policy optimization (PPO)~\cite{schulman2017ppo}
and evaluated across eight observation noise levels (3--30~mm).
As a stretch goal, POMCP-MCTS is also evaluated~\cite{silver2010pomcp} online
planning with the simulator as a generative model.

The main finding is that the benefit of belief augmentation is strongly
noise-regime-dependent.
Belief PPO wins or ties at six of the eight noise levels tested and achieves
its largest gain ($+24\%$) at $\sigmaobs = 10$~mm, a regime where the particle
filter converges quickly yet provides meaningful position smoothing.
At very low noise ($3$~mm), however, plain PPO outperforms by $30\%$:
the raw wrist observation is already accurate, and the filter's initialization
spread induces unnecessary caution in the policy.
This noise-regime dependence is the central empirical contribution of the paper.

\textbf{Departures from the proposal.}
Three aspects of the final system differ from the original proposal, each for
concrete technical reasons.
First, the proposal specified \emph{episodic} observation noise (a different
$\sigmaep$ drawn per episode during training) to produce a single generalist
policy robust to all noise levels simultaneously.
In practice, training with randomized noise destabilized learning: the particle
filter's initialization spread (set proportionally to $\sigmaep$) varied
wildly across episodes, preventing the policy from learning a consistent
relationship between $\bm{\sigma}$ and appropriate caution.
Training at a fixed $\sigmaep$ and evaluate the trained policy across
all noise levels, which cleanly isolates the noise-regime structure.
Second, the POMCP stretch goal was proposed with a \emph{learned} world model
(trained on rollouts from the belief PPO policy).
Learning a faithful transition model for a MuJoCo contact-rich environment
within project scope is impractical; the simulator is used directly
as the generative model, which is the standard approach in model-based POMDP
solvers and does not change the algorithmic comparison.
Third, the proposal described explicit \emph{active perception}, rewarding
the agent for moving to reduce uncertainty (decrease $\bm{\sigma}$) as a
separate reward term.
This was implemented and evaluated; it failed because the
information-gathering reward conflicted with the approach reward, causing the
policy to wander away from the target to reduce occlusion rather than
approaching to grasp.
Removing the explicit reward and relying solely on the task reward produced
a policy that still learns implicit uncertainty-conditioned behavior, as
evidenced by the noise-sweep crossover.

% -- II. Background and Related Work ---------------------------------------
\section{Background and Related Work}
\label{sec:background}

\textbf{POMDPs for robotic manipulation.}
The partially observable Markov decision process (POMDP) framework
\cite{kaelbling1998pomdp} provides the formal foundation for decision-making
under state uncertainty.
In manipulation, partial observability arises from sensor noise, occlusion,
and contact uncertainty.
Koval et al.~\cite{koval2015manifold} show that contact interactions can be
modeled as belief-space transitions, enabling planning over distributions of
object poses rather than single estimates.
This work addresses the simpler but practically important case of visual
observation noise and occlusion, without the added complexity of contact-induced
uncertainty.

\textbf{Belief-state representations.}
Maintaining a belief over the hidden state is the key challenge in POMDP
solving.
The bootstrap particle filter~\cite{gordon1993novel} represents belief as a
weighted set of samples and supports arbitrary non-Gaussian distributions,
making it well-suited to manipulation tasks where occlusion can render
observations completely absent.
Kalman-family filters (KF, EKF, UKF) are computationally cheaper but assume
unimodal Gaussian beliefs; this assumption breaks when the target has been
occluded long enough that the true position could lie anywhere in a region,
creating a diffuse, multi-modal belief that a single Gaussian cannot represent.
In this setting, distractors can occlude the target for $50$+ consecutive steps,
making the particle filter the appropriate choice.

\textbf{Deep RL for manipulation.}
Proximal policy optimization~\cite{schulman2017ppo} is the de facto standard
for continuous-action robotic RL due to its stability and ease of tuning.
Reward shaping is known to be critical and treacherous: Ng et al.~\cite{ng1999potential}
show that only potential-based shaping preserves the optimal policy, while
arbitrary dense rewards can introduce spurious local optima.
Five distinct classes of reward-hacking exploit were discovered during development
(per-step proximity farming, per-step gripper clamping, grasp-attempt spamming,
fine-positioning hovering, and information-gathering conflicts), each illustrating
a different failure mode of ad hoc shaping.

\textbf{Online POMDP planning.}
Silver and Veness~\cite{silver2010pomcp} introduced POMCP, which uses
Monte Carlo tree search with particle filters to plan in large POMDPs without
an explicit model.
Sunberg and Kochenderfer~\cite{sunberg2018pomcpow} extend this to continuous
observation spaces via double progressive widening.
The MCTS implementation incorporates observation progressive widening from
POMCPOW, tree reuse across steps (ABT-style), and progressive bias initialization
inspired by AlphaGo~\cite{silver2016alphago}.

% -- III. Problem Formulation ----------------------------------------------
\section{Problem Formulation}
\label{sec:problem}

\subsection{POMDP Definition}

The task is formalized as a POMDP
$(\States, \Actions, \Obs, \Trans, \ObsFn, \Rew, \Disc)$.

\textbf{State space $\States$.}
The full (hidden) system state is:
\begin{equation}
  s = \bigl(\mathbf{q},\; \block,\; \mathbf{p}_{d_1},\; \mathbf{p}_{d_2},\; h\bigr),
\end{equation}
where:
\begin{itemize}
  \item $\mathbf{q} = [q_1, q_2, q_3, q_4, q_5, g] \in \mathbb{R}^6$ are the
        five joint angles (rad) and the binary gripper state $g \in \{0,1\}$.
  \item $\block = (x_b, y_b, \theta_b) \in \mathbb{R}^3$ is the pose of the
        red target block in the world frame.
  \item $\mathbf{p}_{d_1}, \mathbf{p}_{d_2} \in \mathbb{R}^3$ are the poses
        $(x, y, \theta)$ of the two blue distractor blocks.
  \item $h \in \{0, 1\}$ is a binary flag indicating whether the target block
        is currently held by the gripper.
\end{itemize}
The arm configuration $\mathbf{q}$ is fully observable (joint encoders).
All block poses and $h$ are \emph{hidden}; the agent never directly observes
the true positions of the target or distractors.

\textbf{Action space $\Actions$.}
Actions are continuous six-dimensional vectors $a \in [-1, 1]^6$:
\begin{equation}
  a = [\Delta q_1,\; \Delta q_2,\; \Delta q_3,\; \Delta q_4,\; \Delta q_5,\; g_{\text{cmd}}],
\end{equation}
where $\Delta q_i$ are joint velocity commands scaled to $\pm 0.05$~rad/step,
and $g_{\text{cmd}} \in [-1, 1]$ is the gripper command.
Grasping is implemented as auto-grasp: when the end-effector enters a sphere
of radius $r_{\text{grasp}} = 10$~mm around any block, that block is
automatically attached to the gripper via a position constraint.
The agent therefore controls approach and positioning rather than the
grasp trigger directly.

\textbf{Observation space $\Obs$.}
At each step the agent receives a partial observation
$o_t \in \mathbb{R}^{15}$ (plain PPO) or $\mathbb{R}^{18}$ (belief PPO):
\begin{itemize}
  \item $\mathbf{q} \in \mathbb{R}^6$: current joint angles (fully observed).
  \item $\hat{\block}_t \in \mathbb{R}^3$: wrist camera estimate of the target
        position $(x, y, \theta)$, corrupted by distance-dependent Gaussian
        noise (plain PPO) or replaced by the particle filter mean $\bm{\mu}_t$
        (belief PPO).
  \item $\bm{\sigma}_t \in \mathbb{R}^3$: particle filter uncertainty
        $(\sigma_x, \sigma_y, \sigma_\theta)$ (belief PPO only).
  \item $\ee \in \mathbb{R}^3$: end-effector Cartesian position
        $(x, y, z)$ from forward kinematics.
  \item $\goal \in \mathbb{R}^2$: fixed goal position $(x, y)$ in the
        world frame.
  \item $h \in \{0, 1\}$: holding flag (1 if target is grasped).
\end{itemize}

\textbf{Transition model $\Trans$.}
Because the state space is continuous, the transition is not a matrix but
a physics-defined mapping $s_{t+1} = f(s_t, a_t)$ implemented by the MuJoCo
engine (\texttt{mj\_step}, 10 substeps per action, $\Delta t = 0.002$~s).
Joint angles update deterministically from velocity commands.
Blocks are free rigid bodies under gravity; they remain stationary unless
in contact with the gripper.
When $d_t < r_{\text{grasp}}$, a position constraint rigidly attaches the
nearest block to the end-effector with a fixed vertical offset; releasing the
gripper drops the constraint and the block falls under gravity.

\textbf{Observation function $\ObsFn$.}
The wrist camera observation of the target is:
\begin{equation}
  P(o_t \mid s_t) =
  \begin{cases}
    \mathcal{N}\!\left(\block^{xy\theta},\; \sigmaobs(d_t)^2 \mathbf{I}\right)
      & \text{target visible} \\
    \delta\!\left(o_t - o_{\tau}\right)
      & \text{target occluded}
  \end{cases}
  \label{eq:obs_fn}
\end{equation}
where $\sigmaobs(d_t) = \sigmaep \cdot (d_t / d_{\text{ref}})$ is the
distance-dependent noise standard deviation, $\sigmaep$ is fixed per training
run, $d_{\text{ref}} = 0.15$~m, and $o_{\tau}$ is the last visible reading
replayed during occlusion.

\textbf{Reward $\Rew$ and discount $\Disc$.}
The reward function is detailed in Table~\ref{tab:reward}.
The discount factor is $\Disc = 0.99$ with a maximum episode horizon of
$T = 200$ steps, after which the episode terminates as a timeout failure.

\subsection{Sources of Partial Observability}

Three distinct uncertainty sources make this task a POMDP rather than an MDP.
Each affects the observation function $\ObsFn$ differently and requires a
different handling strategy.

\textbf{1. Distance-dependent sensor noise.}
The wrist camera returns a Gaussian-corrupted estimate of the target's
$(x, y, \theta)$ position with standard deviation
$\sigmaobs(d_t) = \sigmaep \cdot (d_t / d_{\text{ref}})$,
where $d_t$ is the current end-effector-to-target distance and
$d_{\text{ref}} = 0.15$~m.
Crucially, this uncertainty is \emph{reducible}: as the agent approaches the
target, $d_t$ decreases and the observation becomes more accurate.
This creates an implicit information-gathering incentive: the agent must
move closer to localize reliably before grasping, without any explicit
active perception reward.

\textbf{2. Occlusion-induced observation dropout.}
When a distractor block passes between the wrist camera and the target,
the observation is entirely absent for that step.
The agent receives no update to $Z$ during occlusion; it must rely on
its last known estimate or, in the belief PPO case, the particle filter's
prediction step (which correctly grows uncertainty during the gap).
Per-step occlusion rates reach $48\%$ during early training exploration,
meaning the agent frequently acts on stale information for extended periods.
Unlike sensor noise, occlusion-induced uncertainty is \emph{not reducible
by approaching}; it is resolved only by repositioning to clear the
distractor from the camera's line of sight.

\textbf{3. Distractor identity ambiguity.}
Both the target and distractor block poses are hidden states in $s$;
the wrist camera observation covers only the target block and gives no
direct information about distractor positions.
Because the auto-grasp threshold ($r_{\text{grasp}} = 10$~mm) is comparable
to the mean inter-block spacing ($\sim\!17$~mm), a position estimate corrupted
by noise or occlusion can place the agent within grasp range of a distractor
rather than the target.
A wrong grasp terminates the episode with a $-30$ penalty.
This uncertainty source is a direct consequence of the first two: accurate,
unoccluded localization eliminates wrong-grasp risk; degraded localization
makes it non-negligible.

Taken together, these three sources create a genuine exploration--exploitation
tension: the agent must decide whether to spend steps reducing uncertainty
(approaching for a better sensor reading, or repositioning to clear an
occlusion) or to act on its current, imperfect belief and risk a failed grasp.

\subsection{Reward Function}

The reward function uses potential-based shaping~\cite{ng1999potential} for
approach and carry, one-time milestone bonuses to avoid per-step exploitation,
and large terminal rewards to ensure the overall objective dominates.
The full breakdown is given in Table~\ref{tab:reward}.

\begin{table}[!t]
  \caption{Reward function components. All distance thresholds are in mm.}
  \label{tab:reward}
  \centering
  \small
  \begin{tabular}{l c p{3.0cm}}
    \toprule
    Component & Value & Condition \\
    \midrule
    Step cost      & $-1.0$                  & Every step \\
    XY approach    & $+3.0 \!\cdot\! \Delta d_{xy}$ & Per step, potential-based \\
    Proximity bonus & $+2.0\!(1 - d/25\text{mm})$ & $d_{xy} < 25$~mm \\
    Reach block    & $+5.0$ (once)           & $d_{xy} < 15$~mm \\
    Grasp success  & $+10.0$                 & Auto-grasp, target \\
    Grasp fail     & $-2.0$                  & Auto-grasp misses \\
    Distractor grasp & $-30.0$ + terminate   & Wrong block \\
    Carry shaping  & $+5.0 \!\cdot\! \Delta d_{\text{goal}}$ & While holding \\
    Near goal      & $+5.0$ (once)           & Holding, $d_{\text{goal}} < 30$~mm \\
    Perfect place  & $+50.0$                 & $d_{\text{goal}} < 10$~mm \\
    Precise place  & $+30.0$                 & $d_{\text{goal}} < 20$~mm \\
    Close place    & $+10.0$                 & $d_{\text{goal}} < 40$~mm \\
    Miss           & $-10.0$                 & Release outside goal \\
    \bottomrule
  \end{tabular}
\end{table}

\subsection{Simulation Environment}

The SO-ARM101 is a 5-DOF desktop arm with a parallel-jaw gripper.
The MuJoCo scene contains a red $2{\times}4$ LEGO brick (the target,
$32{\times}16$~mm footprint), two blue $2{\times}2$ LEGO bricks (distractors,
$40{\times}20$~mm footprint, $\sim\!25\%$ larger to exaggerate occlusion), and a
fixed yellow goal marker approximately $15$~cm from the home configuration.
Distractors are spawned uniformly within $25$~mm of the target at each episode
reset, producing an inter-block spacing of $17$~mm on average and ensuring
that $\sim\!34\%$ of distractors are within the grasp threshold of the target.

Per-step occlusion rates (measured across $86\%$ of episodes) sit near $16\%$
under a trained policy rollout, rising to $48\%$ under random exploration.
This ensures that belief tracking provides meaningful benefit without making
the task unsolvable.

% -- IV. Solution Approach -------------------------------------------------
\section{Solution Approach}
\label{sec:approach}

\begin{figure*}[!t]
  \centering
  \includegraphics[width=0.88\textwidth]{figs/system_pipeline.pdf}
  \caption{Architecture comparison.
           (a)~Plain PPO exposes noisy (or stale, during occlusion) sensor
           readings as a 15-D vector directly to the policy.
           (b)~Belief-augmented PPO routes each observation through a
           300-particle bootstrap filter; the filter summary
           $(\bm{\mu}_b, \bm{\sigma}_b)$ augments the 18-D policy input so
           the agent can reason about its own localization uncertainty.}
  \label{fig:pipeline}
\end{figure*}

\subsection{Plain PPO (Baseline)}

The plain PPO agent treats the POMDP as a standard MDP by using the stale
sensor model: during occlusion, the policy receives the cached last-visible
wrist observation $o_{\tau}$ as if it were the current reading.
The 15-dimensional policy observation is:
\begin{equation}
  \phi_{\text{plain}} =
  [\,\mathbf{q}_{(6)},\; \hat{\block}^{xy\theta}_{(3)},\;
   \ee{}_{(3)},\; \goal{}_{(2)},\; h_{(1)}\,],
  \label{eq:obs_plain}
\end{equation}
where $\hat{\block}$ is the stale/noisy wrist reading,
$\ee$ is the end-effector Cartesian position (from forward kinematics),
and $\goal$ is the fixed goal position in the world frame.
The policy is an MLP with two hidden layers of 256 units each (ReLU
activations, linear output), matching the SB3 PPO default architecture.

Training uses the SB3 PPO implementation~\cite{raffin2021sb3} with
$n_{\text{env}} = 3$ parallel environments, $n_{\text{steps}} = 2048$,
entropy coefficient $0.05$, $\gamma = 0.99$, clip ratio $0.2$,
and VecNormalize for both observations and rewards.
The plain agent is trained for $5\text{M}$ environment steps; the best checkpoint
is saved by an evaluation callback ($50$ evaluation episodes every $10\text{k}$
steps).

\subsection{Belief-Augmented PPO}

\subsubsection{Particle Filter Belief Tracker}

A bootstrap particle filter is maintained~\cite{gordon1993novel} over the target
block's hidden pose $\block = (x, y, \theta)$.
The filter state is $N = 300$ weighted particles
$\{(x^i, w^i)\}_{i=1}^{N}$ updated at each step as follows.

\textbf{Predict.}
Each particle is perturbed with Gaussian process noise to account for residual
uncertainty in the block's rest position:
\begin{equation}
  x^i_{t|t-1} = x^i_{t-1} + \epsilon,\quad
  \epsilon \sim \mathcal{N}(\mathbf{0},\;\text{diag}(\sigma_p^2, \sigma_p^2, \sigma_{p\theta}^2)),
\end{equation}
with $\sigma_p = 1$~mm/step and $\sigma_{p\theta} = 0.01$~rad/step.
The small process noise reflects that the blocks are stationary objects
that do not move spontaneously.

\textbf{Update.}
When a wrist observation $o_t$ is available (target not occluded), particle
weights are updated by the Gaussian likelihood:
\begin{equation}
  w^i \;\propto\; w^i \cdot \mathcal{N}\!\left(o_t^{xy};\; x^i_{xy},\;
  \sigma_{\text{lik}}^2 \mathbf{I}\right),
  \label{eq:pf_update}
\end{equation}
where $\sigma_{\text{lik}} = 1.5\,\sigmaobs(d_t)$.
The widened likelihood ($1.5\times$) prevents particle impoverishment:
using the raw sensor sigma collapses the effective sample size $N_{\text{eff}}
= 1/\sum_i (w^i)^2$ to near $1$ each step, degenerating the filter into
a noisy running average.
During occlusion, the update step is skipped entirely; particles spread
freely via the predict step, correctly growing uncertainty without
incorporating spurious information.

\textbf{Resample.}
Systematic resampling is applied when $N_{\text{eff}} < N/2$.
After resampling, $5\%$ of particles are re-injected near the last visible
observation to prevent filter degeneracy during extended occlusion.

\textbf{Belief summary.}
The filter's output is compressed to
$\pf_t = (\bm{\mu}_t, \bm{\sigma}_t) \in \mathbb{R}^6$:
\begin{equation}
  \bm{\mu}_t = \sum_i w^i x^i,\quad
  \sigma_{t,k} = \sqrt{\sum_i w^i (x^i_k - \mu_{t,k})^2},
\end{equation}
for each dimension $k \in \{x, y, \theta\}$.
The standard deviation $\bm{\sigma}$ summarizes how uncertain the filter is
about the target's position: high $\bm{\sigma}$ indicates recent or extended
occlusion, while low $\bm{\sigma}$ indicates confident localization.

\subsubsection{Policy Observation and Architecture}

The belief-augmented policy receives an 18-dimensional observation:
\begin{equation}
  \phi_{\text{belief}} =
  [\,\mathbf{q}_{(6)},\; \bm{\mu}_{(3)},\; \bm{\sigma}_{(3)},\;
   \ee{}_{(3)},\; \goal{}_{(2)},\; h_{(1)}\,].
  \label{eq:obs_belief}
\end{equation}
The critical difference from plain PPO is the inclusion of $\bm{\sigma}$:
the policy can learn to slow its approach when uncertainty is high
(reducing the risk of triggering auto-grasp on a distractor at the wrong position)
and act more aggressively when the filter is confident.
No explicit information-gathering reward is used; this behavior must emerge
from the task reward alone.

The reward for belief PPO uses the particle filter mean $\bm{\mu}$ in place
of the raw sensor reading when computing distance-based shaping terms.
This ensures that the reward gradient is also filtered, providing a cleaner
learning signal at high noise levels.

All hyperparameters are identical to plain PPO, except the belief agent trains
for $6.4\text{M}$ steps to account for the larger and harder-to-exploit
observation space.
The policy network architecture is the same MLP~[256, 256] as the baseline.

\begin{figure}[!t]
  \centering
  \includegraphics[width=\columnwidth]{figs/belief_evolution.pdf}
  \caption{Belief uncertainty $\sigma_b$ (top) and distance to target
           (bottom) over one representative successful episode.
           The shaded region marks prolonged wrist-camera occlusion
           (steps 0--178): the particle filter's uncertainty grows steadily
           as no measurements arrive.
           Once occlusion clears, the robot approaches and the
           distance-dependent noise model $\sigmaobs \propto d_t$ causes
           $\sigma_b$ to collapse from 73\,mm to $<\!10$\,mm within a few
           steps.
           The dashed line marks the successful auto-grasp at step~188.}
  \label{fig:belief-evo}
\end{figure}

% -- V. Results ------------------------------------------------------------
\section{Results}
\label{sec:results}

\subsection{Training Convergence}

Fig.~\ref{fig:learning} shows training curves for both agents.
Both discover the approach-shaping gradient around $200\text{k}$ steps,
after which reward rises monotonically until plateauing.
Plain PPO reaches its best checkpoint at step $3.74\text{M}$; belief PPO plateaus later, reaching
its best checkpoint at $6.37\text{M}$ steps.
The slower convergence of belief PPO is expected: the 18-dimensional
observation contains the $\bm{\sigma}$ uncertainty signal, whose relevance
to the task reward requires more experience to exploit.
Neither agent exhibits catastrophic forgetting after the best checkpoint,
indicating stable training.

\begin{figure}[!t]
  \centering
  \subfloat[Mean episode reward]{%
    \includegraphics[width=0.48\columnwidth]{figs/learning_curves.pdf}%
    \label{fig:reward_curve}%
  }
  \hfill
  \subfloat[Task success rate]{%
    \includegraphics[width=0.48\columnwidth]{figs/success_curve.pdf}%
    \label{fig:success_curve}%
  }
  \caption{Training curves (15-step moving average).
           Success rate: fraction of evaluation episodes completing in
           $<200$ steps. Belief PPO converges $\sim1.6\text{M}$ steps
           later but reaches a comparable peak.}
  \label{fig:learning}
\end{figure}

\subsection{Noise Sweep: Main Result}

Both best-model checkpoints are evaluated across eight noise levels
($\sigmaep \in \{3, 5, 8, 10, 13, 16, 20, 30\}$~mm) with 50 evaluation
episodes each at fixed seed 42.
Results are shown in Fig.~\ref{fig:noise_sweep} and Table~\ref{tab:noise_sweep}.

\begin{figure}[!t]
  \centering
  \includegraphics[width=\columnwidth]{figs/noise_sweep.pdf}
  \caption{Task success rate vs.\ $\sigmaobs$.
           Shaded regions show which method leads.
           Belief PPO outperforms plain PPO across the moderate-noise
           regime ($5$--$20$~mm); plain PPO dominates at very low
           noise ($3$~mm).}
  \label{fig:noise_sweep}
\end{figure}

\begin{table}[!t]
  \caption{Noise sweep results (50 episodes per level, seed 42).
           \textbf{Bold} = better. $\Delta$ = Belief $-$ Plain.}
  \label{tab:noise_sweep}
  \centering
  \begin{tabular}{c cc c}
    \toprule
    $\sigmaobs$ (mm) & Plain (\%)$\uparrow$ & Belief (\%)$\uparrow$ & $\Delta$~(\%) \\
    \midrule
    3  & \textbf{66} & 36 & $-30$ \\
    5  & 46 & \textbf{52} & $+6$ \\
    8  & 58 & \textbf{60} & $+2$ \\
    10 & 34 & \textbf{58} & $\mathbf{+24}$ \\
    13 & \textbf{48} & \textbf{48} & $0$ \\
    16 & 44 & \textbf{48} & $+4$ \\
    20 & 40 & \textbf{46} & $+6$ \\
    30 & \textbf{42} & 32 & $-10$ \\
    \bottomrule
  \end{tabular}
\end{table}

The results reveal a clear noise-regime structure.
At $\sigmaep = 3$~mm, plain PPO achieves $66\%$ success with mean episode
length $111$ steps, while belief PPO manages only $36\%$ at $196$ steps.
The plain agent acts decisively because its wrist observations are already
accurate ($\sim\!3$~mm error); the belief agent is overly cautious because
its particle filter is initialized with spread $\sigmaep$ at episode start,
leading the policy to perceive uncertainty even when none is warranted.

In the moderate regime ($\sigmaep = 5$--$20$~mm), belief PPO wins or ties at
every level, with the largest advantage at $10$~mm ($+24\%$, belief $58\%$
vs.\ plain $34\%$).
At this noise level, the particle filter converges within a few visible
observations (typically $<$5 steps) to a mean error of $\sim\!4$~mm,
well below the grasp threshold, while the plain agent's stale observation
carries $10$~mm error, enough to cause localization failures during occlusion.

At $\sigmaep = 30$~mm (beyond the training distribution), both methods
degrade: plain PPO to $42\%$ (stale readings too inaccurate), belief PPO
to $32\%$ (particle impoverishment: the likelihood function is so flat at
$\sigma_{\text{lik}} = 45$~mm that particle weights barely update,
preventing convergence).

The $\sigmaep \approx 8$~mm crossover is theoretically consistent with the
signal-to-noise structure of the bootstrap filter: below it, the raw
observation is already close to optimal and the $O(N)$ filter overhead
is unjustified; above it, filtering reduces position error enough to
consistently reach the $10$~mm grasp sphere before timeout.

\subsection{Failure Mode Analysis}

Across all noise levels, task success rate and grasp rate are nearly
identical ($\text{success}\% \approx \text{grasped}\%$), with zero
wrong-block grasps recorded for either agent.
This reveals that the bottleneck is \emph{localization and approach},
not manipulation skill: once the agent grasps the target block, it
reliably carries and places it (placement phase never fails).
All failures ($\sim\!60\%$ of episodes at mid-to-high noise) are
pre-grasp timeouts in which the agent could not navigate to within $10$~mm
of the target within the $200$-step horizon.

The absence of distractor confusion (zero wrong grasps) is notable given
that $34\%$ of distractor spawns place a distractor within the grasp
threshold of the target.
Both agents learned to approach the target's estimated position accurately
enough that auto-grasp always activates on the correct block.
This is consistent with the high placement success rate conditional on
grasping: both agents learned precise approach behavior; their failure
mode is navigation under uncertainty, not discrimination.

\subsection{POMCP Stretch Goal}

As a stretch goal, POMCP-MCTS is evaluated with the direct MuJoCo simulator
as the generative model.
The planner uses UCB1 tree search with 200 simulations per step, a heuristic
rollout policy, observation progressive widening ($k\!=\!2$, $\alpha\!=\!0.5$),
ABT-style tree reuse across steps, and progressive bias initialization
(heuristic action seeded at $Q\!=\!0$; all others at $Q\!=\!-40$).
The carry-over tree is reset every 40 steps to prevent Q-value corruption
on stuck episodes.
Results on the same 20-episode evaluation (Table~\ref{tab:pomcp}) show
$60\%$ success at a mean planning cost of $28.1$~s/step and 143 mean steps
per episode.

\begin{table}[!t]
  \caption{POMCP-MCTS vs.\ trained policies. PPO columns show results at the
           training noise level ($\sigmaep = 5$~mm). POMCP runs at fixed
           $\sigmaep = 5$~mm (training noise of Belief PPO).}
  \label{tab:pomcp}
  \centering
  \begin{tabular}{l ccc}
    \toprule
    Method & Success (\%)$\uparrow$ & Mean Steps$\downarrow$ & Plan time \\
    \midrule
    Plain PPO   & 46 & $\sim$168 & instant \\
    Belief PPO  & \textbf{52} & $\sim$150 & instant \\
    POMCP-MCTS  & 60 & 143       & 28.1~s/step \\
    \bottomrule
  \end{tabular}
\end{table}

POMCP-MCTS shows a bimodal outcome distribution: successful episodes
complete in 16--86 steps (median $\sim\!40$), while the 8 failures all
hit the 300-step timeout.
This brittleness reflects a known limitation of tree search under a
small simulation budget: with 200 simulations per step, Q-value estimates
carry a standard error of $\sim\!12$ reward units, insufficient to reliably
distinguish between actions whose true values differ by $5$--$15$ reward.
The 300-step timeout failures correspond to hard initial configurations
where the heuristic rollout policy is itself insufficient, and POMCP inherits
this weakness through progressive bias initialization.
Despite the higher nominal success rate, the $28.1$~s/step planning cost
makes POMCP impractical for real-time deployment; trained policies achieve
comparable or better success at negligible inference cost.

% -- VI. Discussion and Conclusion -----------------------------------------
\section{Discussion and Conclusion}
\label{sec:conclusion}

These results show that particle filter belief augmentation provides a meaningful
performance advantage when observation noise is in a regime where filtering is
effective, specifically when $\sigmaep$ is large enough that the raw sensor
reading carries significant error, but not so large that the filter itself
cannot converge.
The crossover at $\sim\!8$~mm and the degradation at $30$~mm are both
theoretically motivated consequences of the bootstrap filter's convergence
characteristics.
The practical implication is that belief augmentation is not universally
beneficial: in low-noise settings, the cost of filter initialization and
the induced caution in the policy can outweigh the smoothing benefit.
Practitioners should choose between a stale-sensor baseline and a
belief-augmented agent based on the expected operating noise level.

On reward engineering, developing a stable reward function for multi-phase
manipulation required iterating through 18 training runs and discovering five
distinct classes of reward hacking.
Per-step state-occupancy rewards (``be near the block'', ``keep gripper closed
near block'', ``be near goal while holding'') were universally exploited:
the agent learned to enter the rewarded state and remain there indefinitely,
collecting reward without task progress.
The fix in every case was to convert the reward to either a one-time milestone
(fires once per episode on the first entry into the state) or a
potential-based shaping term (rewards only the improvement in a scalar function,
not the absolute value).
These are precisely the conditions under which Ng et al.~\cite{ng1999potential}
guarantee that the optimal policy is preserved.
The grasp-attempt spamming exploit in particular (where the agent
discovered it could repeatedly attempt grasps from an incorrect height and
collect positive rewards for each attempt) required four separate design
iterations to resolve robustly.

Regarding limitations and future work, the current environment uses a $2$D
grasp check (XY proximity only) rather than a full $3$D check, a pragmatic
choice after six failed training runs
with the physically correct $3$D threshold.
The issue was the lack of a gradient for the Z descent phase: with a $3$D
grasp condition, the agent learns to approach in XY but finds no incentive
to descend to table height, since the $3$D proximity reward is zero when
hovering directly above the block.
Future work could address this with a curriculum that first rewards XY
approach, then adds a Z descent milestone only after XY convergence is
stable.
Additionally, real robot evaluation via the ROS2 integration
(\texttt{policy\_runner.py}) was begun but not completed due to time constraints;
sim-to-real transfer remains an open question for this platform.
The POMCP-MCTS results highlight that a larger simulation budget (${\sim}2000$
simulations/step) would be needed for reliable Q-value convergence; future work
could explore learned rollout policies to improve planning efficiency.

This work demonstrated that belief-augmented reinforcement learning with a bootstrap
particle filter improves pick-and-place success under observation noise and
occlusion in the moderate-noise regime ($\sigmaep = 5$--$20$~mm), with a peak
advantage of $+24\%$ at $\sigmaep = 10$~mm.
The advantage reverses at very low noise, where raw observations are already
accurate, and at very high noise, where the particle filter degrades.
These findings provide a practical characterization of when belief
augmentation is worth its added complexity, and the failure-mode analysis
identifies localization (not manipulation) as the primary performance
bottleneck in the POMDP pick-and-place setting.

% -- VII. Contributions and Release ----------------------------------------
\section*{Contributions and Release}

The project was designed and built the MuJoCo simulation environment from scratch
(\texttt{lego\_pick\_env.py}), including the reward function, occlusion model,
constraint-based grasping, and the VecNormalize training pipeline.
The bootstrap particle filter (\texttt{particle\_filter.py}) and the
belief-augmented observation wrapper were implemented independently.
The POMCP direct planner and MCTS variant (including tree reuse, observation
progressive widening, periodic reset, and progressive bias initialization)
were implemented from scratch (\texttt{train\_pomcp.py},
\texttt{train\_pomcp\_mcts.py}).
The PPO learning algorithm itself is from Stable-Baselines3~\cite{raffin2021sb3};
the MuJoCo physics engine and SO-ARM101 robot assets are third-party.

The author grant permission for this report to be posted publicly.

% -- References ------------------------------------------------------------
\begin{thebibliography}{99}

\bibitem{kaelbling1998pomdp}
L.~P. Kaelbling, M.~L. Littman, and A.~R. Cassandra,
``Planning and acting in partially observable stochastic domains,''
\emph{Artificial Intelligence}, vol.~101, no.~1--2, pp.~99--134, 1998.

\bibitem{koval2015manifold}
M.~C. Koval, N.~S. Pollard, and S.~S. Srinivasa,
``Pre- and post-contact policy decomposition for planar contact manipulation
under uncertainty,''
\emph{Int'l J. Robotics Research}, vol.~35, no.~1--3, pp.~244--264, 2016.

\bibitem{gordon1993novel}
N.~J. Gordon, D.~J. Salmond, and A.~F.~M. Smith,
``Novel approach to nonlinear/non-Gaussian Bayesian state estimation,''
\emph{IEE Proc.\ Radar and Signal Processing}, vol.~140, no.~2,
pp.~107--113, 1993.

\bibitem{schulman2017ppo}
J.~Schulman, F.~Wolski, P.~Dhariwal, A.~Radford, and O.~Klimov,
``Proximal policy optimization algorithms,''
\emph{arXiv:1707.06347}, 2017.

\bibitem{ng1999potential}
A.~Y. Ng, D.~Harada, and S.~Russell,
``Policy invariance under reward transformations: Theory and application
to reward shaping,''
in \emph{Proc.\ Int'l Conf.\ on Machine Learning (ICML)}, 1999.

\bibitem{silver2010pomcp}
D.~Silver and J.~Veness,
``Monte-Carlo planning in large POMDPs,''
in \emph{Advances in Neural Information Processing Systems (NeurIPS)}, 2010.

\bibitem{sunberg2018pomcpow}
Z.~Sunberg and M.~Kochenderfer,
``Online algorithms for POMDPs with continuous state, action, and
observation spaces,''
in \emph{Proc.\ Int'l Conf.\ on Automated Planning and Scheduling (ICAPS)},
2018.

\bibitem{silver2016alphago}
D.~Silver et al.,
``Mastering the game of Go with deep neural networks and tree search,''
\emph{Nature}, vol.~529, pp.~484--489, 2016.

\bibitem{raffin2021sb3}
A.~Raffin, A.~Hill, A.~Gleave, A.~Kanervisto, M.~Ernestus, and N.~Dormann,
``Stable-Baselines3: Reliable reinforcement learning implementations,''
\emph{J.\ Machine Learning Research}, vol.~22, no.~268, pp.~1--8, 2021.

\end{thebibliography}

\end{document}
